\documentclass[11pt]{article}
\usepackage[margin=1in]{geometry}
\usepackage[T1]{fontenc}
\usepackage{lmodern,microtype,hyperref,enumitem,amsmath,amssymb,xcolor,booktabs}
\usepackage[numbers,sort&compress]{natbib}
\usepackage{url}
\hypersetup{colorlinks=true,urlcolor=blue,citecolor=blue,linkcolor=blue}

\title{Q03: Reasoning in Language Models --- Real Capacity, Linguistic Artifact, or Useful Illusion?}
\author{DLS 2026 --- module author: Kukla}
\date{October 2026}

\begin{document}
\maketitle

\section*{Learning objectives}
\begin{itemize}[leftmargin=*]
  \item Distinguish the competing claims about what ``reasoning'' means in large language models (LLMs), and articulate the empirical signatures that would favor each.
  \item Analyze the mechanisms behind chain-of-thought (CoT) prompting, including when intermediate tokens causally support an answer versus when they rationalize a conclusion reached otherwise.
  \item Evaluate the methodological pitfalls---contamination, surface heuristics, and unfaithful explanations---that make the ``is it real?'' question hard to settle, and apply diagnostic probes to a concrete case.
\end{itemize}

\section{Framing the question}
The student's question is really four hypotheses dressed as one. \emph{Real reasoning} claims the model performs something functionally analogous to inference: composing intermediate conclusions to reach novel answers. \emph{A property of language} claims that reasoning-like behavior is latent in the statistical structure of text, so a good language model inherits it for free. \emph{A reflection} claims the model mirrors reasoning present in its training corpus without instantiating the process. \emph{An illusion} claims the appearance of reasoning is produced by shallow pattern-matching that happens to correlate with correct answers on familiar distributions.

These are not mutually exclusive. A defensible modern view is that LLMs implement \emph{some} genuine computation that deserves the label ``inference,'' that this computation is heavily shaped by the statistics of language, that it frequently reduces to retrieval of memorized patterns, and that our measurements routinely overstate it. The useful graduate-level move is to stop asking the binary question and instead ask: \emph{under what mechanism, measured how, does behavior X count as reasoning?} The sections below supply the mechanistic and empirical vocabulary for doing so.

\section{Mechanism 1: Chain-of-thought as externalized computation}
Chain-of-thought prompting elicits intermediate reasoning steps before a final answer and dramatically improves performance on arithmetic, symbolic, and commonsense tasks for sufficiently large models \citep{wei2022chain}. The same effect appears without any exemplars when the model is simply instructed to think step by step \citep{kojima2022large}, and sampling several chains and taking a majority vote (self-consistency) improves accuracy further \citep{wang2023selfconsistency}. These results establish an operational fact: additional generated tokens can function as scratch space that the model's later computation conditions on.

There is a plausible theoretical reason this helps. A single forward pass of a fixed-depth transformer has bounded serial computation; emitting intermediate tokens and re-ingesting them effectively unrolls additional steps, letting the model trade test-time tokens for depth. Under this reading CoT is not a window into a pre-existing thought but an \emph{enabling} mechanism---the computation partly happens \emph{because} the tokens are produced. This is evidence against the pure ``illusion'' view: if the intermediate tokens were decorative, ablating or corrupting them would not change the answer, yet it often does.

Crucially, the benefit is uneven. CoT helps most on tasks with compositional, symbolic structure and much less on tasks answerable by direct recall. This pattern itself is diagnostic: it suggests the mechanism being recruited is sequential composition, which is one reasonable minimal definition of reasoning.

\section{Mechanism 2: Faithfulness---do the stated reasons cause the answer?}
A reasoning trace can be fluent, plausible, and causally irrelevant. \citet{turpin2023language} show that models produce CoT explanations that systematically omit the true drivers of their answers: inserting a biasing feature (e.g., always making option (A) correct in few-shot exemplars) shifts the model's predictions while the generated rationale never mentions the bias and instead confabulates task-relevant justifications. This is \emph{unfaithfulness}: the stated reasoning does not reflect the actual computation.

Unfaithfulness is the strongest support for the ``reflection/rationalization'' hypothesis. The model has learned what good explanations \emph{look} like---a property of language---and emits them post hoc. Yet the same literature complicates the pessimism: faithfulness is measurable and manipulable, not all-or-nothing. Techniques that force the final answer to depend on the trace (e.g., decomposing a question so the answer is computed only from sub-answers) can increase faithfulness. The methodological lesson for this module is that ``is it reasoning?'' must be posed as a causal question about the trace, answered with interventions, not with how convincing the prose reads.

\section{Mechanism 3: Memorization, contamination, and surface heuristics}
A persistent threat to any ``real reasoning'' claim is that benchmark success reflects exposure rather than inference. Web-scale pretraining corpora overlap with evaluation sets, and even partial contamination inflates scores in ways that masquerade as competence. More subtly, models exploit \emph{surface heuristics}: features that correlate with answers in the training distribution but are not the intended reasoning signal. The classic demonstration in NLI is that models learn annotation artifacts---lexical cues such as negation words predicting ``contradiction''---and achieve high accuracy without entailment reasoning \citep{gururangan2018annotation}; such models collapse on adversarial examples where the cue is decorrelated from the label \citep{mccoy2019right}.

For arithmetic and symbolic tasks the analogous test is \emph{robustness to irrelevant perturbation}. \citet{mirzadeh2024gsm} construct GSM-Symbolic, generating variants of grade-school math problems by changing names and numbers and by inserting plausible-but-irrelevant clauses; strong models' accuracy degrades markedly, and performance is sensitive to superficial numerical changes. This is exactly what the ``illusion'' hypothesis predicts for a system matching templates rather than executing a procedure. The honest conclusion is graded: models hold a mixture of genuine procedural competence and brittle pattern-matching, and the proportion depends on task, scale, and how far the probe strays from the training distribution.

\section{Mechanism 4: Emergence, scale, and whether the question is well posed}
Part of the excitement around LLM reasoning came from reports that certain capabilities appear abruptly at scale---seemingly absent in small models and present in large ones \citep{wei2022emergent}. If true in a strong sense, emergence would suggest a qualitative phase change toward reasoning. But \citet{schaeffer2023emergent} argue that many ``emergent'' jumps are artifacts of the metric: discontinuous, all-or-nothing scoring (exact-match accuracy) manufactures sharp curves out of smoothly improving underlying log-likelihoods. Under continuous metrics the transitions are often gradual.

This methodological point reframes the student's question. ``Did reasoning suddenly appear?'' may be the wrong frame; the measured quantity is a thresholded view of a continuous competence. It also warns against treating any single benchmark number as evidence that reasoning is or is not ``real.'' The appearance of a capability is partly a property of our ruler.

\section{Mechanism 5: What interpretability adds}
Behavioral probes bound what a model does; mechanistic interpretability asks how. Work isolating interpretable internal structure---such as induction heads implementing a copy-and-continue pattern that underlies much in-context learning \citep{olsson2022incontext}---shows that specific, generalizable circuits form during training rather than mere lookup tables. In-context learning itself, the substrate on which few-shot ``reasoning'' rides, has been characterized as the model inferring a latent task from the prompt \citep{xie2022explanation}. These findings tilt against the pure ``illusion'' reading: there exist reusable algorithmic components. They do not, however, license the claim that the model reasons about everything; a circuit for pattern continuation is not a general theorem prover.

The practical upshot: when you want to adjudicate a specific claim (``the model solves this by addition, not recall''), the strongest evidence combines (a) out-of-distribution behavioral generalization, (b) causal trace interventions for faithfulness, and (c) where feasible, identification of the internal mechanism.

\section{Synthesis: a graded answer}
So---real, linguistic, reflection, or illusion? The defensible position is a conjunction:
\begin{itemize}[leftmargin=*]
  \item \textbf{Partly real.} CoT and self-consistency recruit genuine sequential computation that changes answers causally, and interpretable circuits exist; this exceeds mere surface mimicry \citep{wei2022chain,wang2023selfconsistency,olsson2022incontext}.
  \item \textbf{Deeply linguistic.} The competence is shaped by, and often reducible to, statistical structure of text, including task inference from context \citep{xie2022explanation}.
  \item \textbf{Often reflective/unfaithful.} Stated rationales frequently do not cause the answer \citep{turpin2023language}.
  \item \textbf{Sometimes illusory.} Robustness tests expose template-matching and artifact exploitation \citep{gururangan2018annotation,mccoy2019right,mirzadeh2024gsm}, and ``emergence'' can be a metric artifact \citep{schaeffer2023emergent}.
\end{itemize}

For your own experiments: never report a reasoning claim from aggregate accuracy alone. Pair it with a perturbation probe (does performance survive irrelevant edits?), a faithfulness intervention (does corrupting the trace change the answer?), and a contamination check. Reasoning in LLMs is best treated as a measured, mechanism-specific, graded property---not a metaphysical yes/no.

\bibliographystyle{plainnat}
\bibliography{bibliography}
\end{document}